% =============================================================================
% Section 10 — Limits of Invariance
% Book I — Ontology of Viability
% Chapter 02 — Invariance
%
% Purpose:
%   To establish the conceptual and ontological limits of invariance
%   and prevent its reduction to identity, persistence, continuity,
%   structure, or the exhaustive determination of an existent.
%
% Conceptual Policy:
%   1. Invariance is a limited mode of preservation.
%   2. Invariance cannot independently establish ontological identity.
%   3. The absence of identified invariants does not establish
%      the absence of identity.
%   4. Formal invariance does not automatically establish
%      ontological invariance.
%   5. No universal invariant is presumed to exist.
%   6. The limits of invariance do not diminish its analytical
%      and ontological significance.
% =============================================================================

\section{Limits of Invariance}
\label{sec:limits-of-invariance}

\subsection{The Need to Establish the Limits of Invariance}

The preceding analysis has established invariance as a
specific mode of preservation under transformation.
It has also distinguished invariance from identity,
difference, structure, composition, continuity, and
the domains within which preservation may be examined.

These distinctions make it necessary to establish the
limits of invariance explicitly.

Without such limits, invariance may be assigned an
explanatory role exceeding its proper scope. The preservation
of a determination may be treated as proof of the persistence
of an existent, while the transformation of that determination
may be interpreted as evidence of the loss of identity.

Both conclusions exceed what invariance alone can establish.

Invariance identifies the preservation of a specified
determination relative to a specified transformation.
It does not independently determine the complete ontological
status of the existent in which that determination occurs.

The purpose of this section is therefore not to diminish
the importance of invariance, but to establish the boundaries
within which its explanatory and analytical use remains valid.

\subsection{Invariance Is Not Identity}

The first and most fundamental limit concerns the distinction
between invariance and identity.

Identity concerns the existent as such. Invariance concerns
the preservation of a particular determination under
transformation.

The preservation of one or more determinations does not
establish that the existent possessing them is numerically
identical across the transformations under consideration.

Distinct existents may possess equivalent or even
indistinguishable invariant determinations. Their
structural, quantitative, qualitative, or compositional
features may coincide without making them one and the
same existent.

Conversely, an existent may undergo transformation in
which previously identified invariant determinations
are altered or cease to be preserved.

Such alteration does not independently establish that
the existent has lost its identity.

Consequently, neither the presence nor the absence of
an invariant is sufficient to determine numerical identity.

Identity must remain an independent ontological question.

\subsection{Invariance Is Not Persistence}

Invariance must also be distinguished from persistence.

Persistence concerns the continued existence of an existent
across a specified temporal or transformational domain.
Invariance concerns the preservation of a particular
determination within such a domain.

A determination may remain invariant across a sequence
of states without establishing that the same existent
persists throughout that sequence.

For example, distinct existents may exhibit the same
quantitative or structural determination. The preservation
of that determination across them does not establish
the persistence of one existent.

Likewise, an existent may persist through transformations
in which particular determinations are altered.

The persistence of an existent therefore cannot be reduced
to the preservation of an invariant.

Any claim concerning persistence requires an independent
account of the relation between the existent and its
successive states or transformations.

\subsection{Invariance Is Not Continuity}

Continuity concerns the connectedness of successive
states or phases. Invariance concerns the preservation
of a specified determination across transformation.

These concepts may coincide in particular cases, but
neither entails the other.

A process may remain continuous while its determinations
change. The absence of an identified invariant does not
necessarily establish a discontinuity.

Conversely, an invariant may be preserved across states
whose ontological connectedness has not been established.

The preservation of a determination therefore cannot
serve as a universal criterion of continuity.

Nor can the transformation or disappearance of an
invariant independently establish that continuity
has ceased.

Continuity must be examined according to the relation
connecting the relevant states or phases.

\subsection{Invariance Is Not Stability or Immutability}

Invariance is also distinct from stability and immutability.

Stability concerns the maintenance of a condition,
configuration, or organization under specified
circumstances or perturbations.

Immutability concerns the absence of change within
the relevant respect.

Invariance, by contrast, concerns preservation under
a specified transformation. Such preservation may
occur even when other determinations change.

An existent may therefore exhibit invariance without
being stable in every respect. It may also undergo
substantial transformation while preserving a particular
determination.

Likewise, the preservation of one invariant does not
establish that the existent is immutable.

The concepts must remain distinct because each addresses
a different relation between preservation and change.

\subsection{The Absence of a Universal Invariant}

No universal invariant should be presumed to exist
across all existents, transformations, and ontological
domains.

An invariant is always identified relative to a particular
determination, transformation, domain, and respect
of preservation.

A determination preserved in one case may not be preserved
in another. A transformation under which a particular
feature remains invariant may differ from a transformation
under which that feature changes.

Consequently, the existence of an invariant in one
existent or domain does not establish its universal
applicability.

Nor should the analysis assume that every existent
must possess at least one invariant under every possible
transformation.

Such a presumption would transform a limited analytical
concept into an unsupported universal ontological principle.

The possibility of invariance must be investigated
in each relevant case rather than imposed in advance.

\subsection{The Limits of Formal Invariance}

Formal invariance is established within a defined
formal system and relative to specified transformations
and preservation criteria.

Its validity may be mathematically or logically rigorous
within that system. However, formal validity alone
does not establish that the same invariant belongs
to an independently existing entity.

A formal representation may preserve a mathematical
relation while the ontological status of the corresponding
relation remains undetermined.

The transition from formal invariance to ontological
invariance therefore requires an account of the relation
between the formal system and the existent being described.

This requirement does not reject formal analysis.
It prevents the unqualified transfer of formal conclusions
into ontological claims.

Formal invariance may provide a precise description
of a possible pattern of preservation, but its ontological
significance must be independently established.

\subsection{The Limits of Observed Preservation}

The identification of an invariant may depend upon
the available descriptions, measurements, or analytical
methods through which a transformation is examined.

An observed determination may appear preserved within
a limited range of observations while its preservation
under a broader transformation domain remains
undetermined.

Likewise, a determination that appears to change under
one mode of description may remain preserved under
another appropriately specified criterion.

The distinction between observed preservation and
ontological preservation is therefore necessary.

Observed preservation provides evidence for an invariance
claim, but the scope of that claim cannot exceed the
conditions under which preservation has been established.

The absence of observed change is not automatically
proof of universal invariance.

Similarly, the failure to identify an invariant does
not establish its ontological absence.

These limitations concern the justification and scope
of invariance claims. They do not imply that invariance
is merely subjective or dependent upon observation.

\subsection{The Limits of Structural and Compositional Invariance}

Structural and compositional invariance provide
important descriptions of preservation, but neither
can serve as an exhaustive account of an existent.

Structural invariance concerns the preservation of
a specified organizational determination.

Compositional invariance concerns the preservation
of a specified determination relating to the constitution
of a composed existent.

Neither form establishes numerical identity by itself.

An existent may undergo structural reorganization
without the mere occurrence of that reorganization
establishing the loss of identity.

Likewise, the replacement or transformation of components
does not independently determine whether the existent
has ceased to be the existent that it is.

Conversely, the preservation of a structure or composition
does not establish that two descriptions or successive
states refer to one numerically identical existent.

The limits of these forms follow from the more general
distinction between the preservation of a determination
and the identity of the existent.

\subsection{The Limits of Cross-Domain Inference}

An invariant established within one domain cannot
automatically be transferred to another.

The preservation of a structural determination does
not establish the preservation of a corresponding
material determination.

Likewise, the preservation of a relational determination
does not automatically establish the preservation of
a determination belonging to First Determination.

Such transfers require an explicit account of the
relation between the relevant domains and determinations.

The same requirement applies to the four fundamental
domains and the three intermediate hubs.

Their distinction must be preserved when examining
invariance. The existence of a correspondence between
two domains does not establish that their determinations
are ontologically identical.

Cross-domain invariance must therefore be demonstrated
rather than presumed from similarity, analogy, or
formal correspondence.

\subsection{The Scope of Transformation}

An invariance claim is necessarily restricted by the
transformation domain within which preservation is
asserted.

A determination may remain invariant under a limited
class of transformations while changing under another.

The preservation of a determination under one transformation
therefore cannot be generalized to all possible
transformations without additional justification.

The same applies to the range of conditions under
which preservation is established.

An invariant may be conditional, domain-specific,
or restricted to a particular respect.

These restrictions do not invalidate the invariance
claim. They define its actual scope.

A precise account must therefore distinguish:

\begin{itemize}
    \item the transformation under which preservation
          is established;
    \item the transformations outside the established
          domain;
    \item the determination that remains preserved;
    \item the criterion by which preservation is assessed;
    \item the range within which the claim remains valid.
\end{itemize}

The absence of such specification makes an invariance
claim ambiguous and potentially misleading.

\subsection{Invariance Cannot Exhaust the Ontological Constitution of an Existent}

An existent cannot be exhaustively characterized
through the inventory of its invariant determinations.

Such an approach would assume that the ontological
constitution of an existent is reducible to those
determinations that remain preserved under selected
transformations.

This assumption is not justified.

An existent may possess determinations that change,
determinations that remain preserved, and determinations
whose relations to transformation have not yet been
established.

The ontological significance of an existent cannot
therefore be reduced to the set of features classified
as invariant.

Invariance identifies a particular relation between
determination and transformation. It does not provide
a complete account of being.

The analysis of invariance must consequently remain
situated within a broader ontological examination
of the existent.

\subsection{The Limits of Negative Inference}

The limits of invariance are especially important
when negative conclusions are drawn from the failure
to identify preserved determinations.

The absence of an identified invariant does not establish:

\begin{itemize}
    \item the absence of identity;
    \item the loss of numerical identity;
    \item the discontinuity of an existent;
    \item the termination of persistence;
    \item the absence of ontological organization;
    \item the absence of determinations not yet identified.
\end{itemize}

Such conclusions require independent arguments.

The failure to establish invariance may result from
the absence of preservation, the inadequacy of the
specified transformation domain, the selection of
an inappropriate criterion, or the limits of the
available analysis.

These possibilities must not be collapsed into one
ontological conclusion.

A negative invariance claim is therefore limited to
the determination, transformation, domain, and respect
that have actually been examined.

\subsection{The Limits of Positive Inference}

The same caution applies to positive conclusions.

The identification of an invariant does not establish:

\begin{itemize}
    \item the numerical identity of distinct states
          or existents;
    \item the persistence of one existent across every
          transformation;
    \item the continuity of a process;
    \item the immutability of the existent;
    \item the preservation of all other determinations;
    \item the identity of different ontological domains.
\end{itemize}

An invariant establishes preservation only within
its specified scope.

Any conclusion extending beyond that scope requires
additional ontological justification.

The presence of an invariant must therefore be treated
as a determinate finding rather than as a universal
explanation.

\subsection{The Positive Significance of Invariance}

Establishing the limits of invariance does not render
the concept philosophically marginal.

Invariance remains important because it allows the
analysis of preservation within transformation.

It makes possible the distinction between determinations
that remain preserved and those that undergo change.

It also provides a means of examining the conditions
under which particular features, relations, structures,
or compositional arrangements persist.

Without this distinction, the analysis of transformation
would risk treating all change as undifferentiated.

The proper limitation of invariance therefore strengthens
rather than weakens its analytical value.

By refusing to reduce identity, persistence, or continuity
to invariance, the ontological analysis preserves the
distinct significance of each concept.

\subsection{Synthesis}

The limits established in this section may be summarized
as follows:

\begin{itemize}
    \item Invariance is not identity.
    \item Invariance is not persistence.
    \item Invariance is not continuity.
    \item Invariance is not stability or immutability.
    \item No universal invariant is presumed.
    \item Formal invariance is not automatically ontological
          invariance.
    \item Observed preservation does not establish
          unrestricted ontological preservation.
    \item Invariance in one domain does not automatically
          transfer to another.
    \item The absence of an identified invariant does
          not establish the absence or loss of identity.
    \item The presence of an invariant does not establish
          numerical identity.
    \item Invariance cannot exhaust the ontological
          constitution of an existent.
\end{itemize}

Invariance must therefore be understood as a limited
but significant mode of preservation within transformation.

Its validity depends upon the specification of the
determination, transformation, domain, respect, and
criterion involved.

The recognition of these limits prevents invariance
from becoming an unjustified universal principle and
preserves the ontological distinction between the
preservation of determinations and the identity of
the existent.

% =============================================================================
% End of Section 10
% =============================================================================